% !TeX root = main.tex
\section{Experimental Plan and Current Status}

\subsection{Staged Validation Questions}

RSVP의 최종 비교 전에 다음 세 연결고리를 순서대로 검증한다.

\begin{enumerate}
  \item \textbf{기반 지침의 유용성:} 동일한 QMIX backbone에서 실제 LLM 지침이
        무지침 QMIX 및 상태와 무관하게 섞은 지침보다 외부 성능을 높이는가?
  \item \textbf{시간순 예측력:} critic이 아직 학습하지 않은 episode의 predicate
        접미합을 평균 또는 지수이동평균 기준선보다 정확하게 예측하는가?
  \item \textbf{갱신의 인과적 효용:} 낮은 $v_t$ 또는 높은 $S_t$에서 지침을
        갱신한 이후의 추가 학습 성능이 같은 상태에서 지침을 유지한 경우보다 좋은가?
\end{enumerate}

첫 단계가 실패하면 갱신기 파라미터 탐색보다 지시문, 접지, 셰이핑 크기와 학습 통합을
먼저 수정해야 한다. 두 번째 단계가 실패하면 잔여 가치 대신 더 단순한
상태 변화 또는 성과 기반 trigger를 고려한다. 예측력이 있더라도 세 번째 단계가
성립하지 않으면 RSVP의 성능 기여를 주장하지 않는다.

\subsection{Baselines and Evaluation Metrics}

실험은 StarCraft Multi-Agent Challenge(SMAC)~\cite{samvelyan2019smac}의 서로 다른
난이도와 병력 구성을 가진 맵에서 수행하며, QMIX, 고정 주기 LEHCA, RSVP를 비교한다.
기반 지침이 확인된 뒤에는 동일 호출 예산의 고정 주기, 무작위 갱신,
상태 변화 trigger를 추가한다. 모든 주요 비교는 동일 환경 설정, 학습 지평, 탐색
스케줄과 seed를 공유한다.

최종 승률뿐 아니라 전체 및 초기 학습 곡선의 AUC, seed 간 분산, 최초 승리 시점,
외부 보상과 $\lambda_tF_t$의 상대 크기를 보고한다. 시스템 비용은 논리적 refresh 수,
실제 LLM 요청 수, cache hit, 토큰, 지연을 분리한다. Scheduler 내부에서는 critic의
시간순 오차, 학습 가능한 head coverage, $v_t$와 $S_t$의 분포, 조기/fallback 호출
비율 및 환경 국면 전환과의 정렬을 측정한다.

\subsection{Current Status}

현재까지의 결과만으로 RSVP가 고정 주기 또는 QMIX보다 우수하다고 결론 내릴 수 없다.
가장 큰 병목은 RSVP의 감지기보다 기반 LEHCA 지침 채널이 QMIX를 안정적으로 개선하지
못한다는 점이다. 특히 SMAC의 dense damage/kill 보상이 semantic shaping과 중복되어
LEHCA의 상대적 이득을 가렸을 가능성과, 논문이 전제한 희소 보상 조건의 정확한 설정이
충분히 공개되지 않았다는 문제가 있다. 이에 따라 현재는 sparse terminal reward에서
QMIX와 원형 LEHCA를 동일 조건으로 비교하여 기반 채널의 유효성을 우선 검증하고 있다.
이 결과가 확보되기 전에는 성능 우위에 관한 표와 문구를 확정하지 않는다.
